\documentclass[11pt]{article}
\usepackage[margin=1in]{geometry}
\usepackage{amsmath,amssymb,amsthm,booktabs,longtable,array,calc}
\usepackage[hidelinks]{hyperref}
\usepackage[T1]{fontenc}
\usepackage[utf8]{inputenc}
\newtheorem{theorem}{Theorem}[section]
\newtheorem{proposition}[theorem]{Proposition}
\newtheorem{lemma}[theorem]{Lemma}
\newtheorem{corollary}[theorem]{Corollary}
\theoremstyle{definition}
\newtheorem{definition}[theorem]{Definition}
\theoremstyle{remark}
\newtheorem*{remark}{Remark}
\providecommand{\tightlist}{\setlength{\itemsep}{0pt}\setlength{\parskip}{0pt}}
\begin{document}
\appendix
\section{Detailed Boundary Proof and the Eight Physical Singleton States}\label{appendix-f.-detailed-boundary-proof-and-the-eight-physical-singleton-states}

This appendix uses only the \(n=2\) physical map.

\subsection{F.1. Physical projection and crease
census}\label{f.1.-physical-projection-and-crease-census}

The crease-incidence cycle is \[
A_1-H_1-B_1-L_1-B_2-H_2-A_2-U_1-A_1.
\] The four physical labels are \(h_0,h_1,u_0,l_0\). Hence the assigned
map depends only on \(\pi_1(P)=(h_0,h_1,u_0,l_0)\), and each physical
assignment has four formal representatives obtained by arbitrary choices
of the two absent formal bits.

The complete global test consists of the four primitive M/V precedences
and exactly two Butterfly pairs: \((H_1,H_2)\) and \((U_1,L_1)\). These
are the two pairs of opposite edges of the cycle.

\subsection{F.2. Negative product: Hamiltonian uniqueness and
validity}\label{f.2.-negative-product-hamiltonian-uniqueness-and-validity}

Traverse the cycle as \(A_1\to B_1\to B_2\to A_2\to A_1\) and assign
sign \(+1\) to a primitive edge pointing with this traversal and \(-1\)
to an edge pointing against it. Direct checkerboard evaluation gives
traversal signs \[
h_0,\quad -l_0,\quad h_1,\quad -u_0,
\] whose product is \(h_0h_1u_0l_0=\delta_0\).

If \(\delta_0=-1\), exactly one or three edges oppose the traversal. In
either orientation, the three majority edges form a directed Hamiltonian
path \(v_0\prec v_1\prec v_2\prec v_3\) and the remaining edge is the
shortcut \(v_0\prec v_3\). Thus there is exactly one linear extension.
The two opposite edge pairs have endpoint intervals \([0,1]\) with
\([2,3]\) and \([1,2]\) with \([0,3]\), respectively; one pair is
disjoint and the other nested. Hence the unique order passes both
Butterfly tests.

\subsection{F.3. Positive product: direct global
obstruction}\label{f.3.-positive-product-direct-global-obstruction}

If \(\delta_0=+1\), the number of traversal reversals is \(0,2,\) or
\(4\). With \(0\) or \(4\) reversals the primitive precedences form a
directed \(4\)-cycle. With exactly two reversals, the orientation is
acyclic but every linear extension makes one opposite-edge pair
alternate.

If the two reversed edges are adjacent, the poset has one source, one
sink, and two incomparable middle vertices. The two possible middle
orders force intervals \([0,2]\) and \([1,3]\) for one of the
opposite-edge pairs. If the reversed edges are opposite, the poset has
two sources followed by two sinks; whichever internal source/sink order
is chosen, one opposite-edge pair again has intervals \([0,2]\) and
\([1,3]\). Hence no globally valid state exists.

\subsection{F.4. Exact physical state
table}\label{f.4.-exact-physical-state-table}

\begin{longtable}[]{@{}ll@{}}
\toprule\noalign{}
physical projection \((h_0,h_1,u_0,l_0)\) & exact state set \\
\midrule\noalign{}
\endhead
\bottomrule\noalign{}
\endlastfoot
\(MMMV\) & \(\{(A_1,B_1,B_2,A_2)\}\) \\
\(MMVM\) & \(\{(B_2,A_2,A_1,B_1)\}\) \\
\(MVMM\) & \(\{(A_1,A_2,B_2,B_1)\}\) \\
\(MVVV\) & \(\{(A_2,A_1,B_1,B_2)\}\) \\
\(VMMM\) & \(\{(B_2,B_1,A_1,A_2)\}\) \\
\(VMVV\) & \(\{(B_1,B_2,A_2,A_1)\}\) \\
\(VVMV\) & \(\{(B_1,A_1,A_2,B_2)\}\) \\
\(VVVM\) & \(\{(A_2,B_2,B_1,A_1)\}\) \\
\end{longtable}

The other eight physical projections have empty exact state set. The
table is a literal presentation of the analytic Hamiltonian rule, not a
computational premise.

\end{document}
